\section{Every localized tester cluster is a literal strong-dual datum}
\label{sec:r12-capacity-cluster-strong-dual}

Retain exactly the strict original box and background-extension
hypotheses (M2), with arbitrary original symmetric middle-positive
separately closed $S$ of degree $(p,p)$. Put $q=n-p\geq1$,
$d=q+1$, and $C_0=d2^d$. All capacities in this section use the
full original normalized classical psh tester class on the same
whole domain $V$. Let
\[
 \mathcal A(V)=\{a: a\text{ is original classical psh on }V,
                                      \ 0\leq a\leq1\}.
 \tag{CW1}
\]
Fix an original compact smooth scalar $\chi$ with $0\leq\chi\leq M$,
where $M<\infty$. For any row $a_j\in\mathcal A(V)$ put
\[
       \nu_j=\chi\,(Da_j)^q\wedge S.
 \tag{CW2}
\]
These are positive Radon measures, not necessarily global source
measures of potentials in the zero-boundary energy cone. The latter
distinction is essential below.

\begin{proposition}[An unconditional full-tester cluster datum theorem]
\label{prop:r12-capacity-cluster-literal-dual}
Every subsequential weak original-current limit $\nu$ of (CW2)
is a finite positive Radon measure supported in
$H=\operatorname{supp}\chi$, and for every Borel $E\subset V$,
\[
                         \nu(E)\leq M C_S^*(E;V).
 \tag{CW3}
\]
In particular, with $m=\nu(V)$, literal integration on the complete
original Borel energy-difference space defines
\[
 \begin{split}
 \ell_\nu(h)&=-\int_Vh\,d\nu\in\mathbb D_{S,\mathrm s}^*,\\
 \|\ell_\nu\|_*
       &\leq \frac{d}{d-1}
                 m^{1-1/d}(M C_0)^{1/d}=:B_{m,M}.
 \end{split}
 \tag{CW4}
\]
The right side is zero if $m=0$. For every strongly convergent
row $v_j\to v$ in the unchanged full cone $\mathscr C$,
\[
                         \int_V|v_j-v|\,d\nu\longrightarrow0.
 \tag{CW5}
\]
There is no no-loss, common-barrier, actual-representability or
source-compatibility hypothesis in this proposition.
\end{proposition}

\noindent\textbf{Status: \texttt{DIRECT\_COROLLARY}} of
positive weak compactness, the definition of the full tester
capacity and the proved capacity-modulus strong-dual theorem
(CD1)--(CD10). The full arbitrary-$S$ weak-cluster application
is recorded here to isolate its unconditional content.

\begin{proof}
Relative CLN gives one finite bound for the masses of (CW2),
independent of every $a_j$ and of its actual moment presentations.
All measures are supported in the fixed compact $H$. Original
smooth scalar density and positivity identify weak convergence
on original tests with weak Radon convergence on this compact
stage. For instance, approximate a continuous compact test
uniformly by original smooth tests and use the common mass bound
to pass the approximation error. Positivity and the fixed support
therefore pass to the limit, and $m<\infty$.

For every ordinary open $G\subset V$, the open-set Portmanteau
inequality and the original tester definition give
\[
    \nu(G)\leq\liminf_j\nu_j(G)
        \leq M C_S(G;V).
 \tag{CW6}
\]
For a Borel $E\subset G$, $\nu(E)\leq\nu(G)$.
Taking the infimum over all ordinary open neighborhoods $G$
of $E$ proves (CW3). This does not replace $C_S^*$ by raw
$C_S$ on a compact or Borel set. In particular it uses no
instance of the missing compact no-gain assertion (FG1).

The modulus $F(s)=Ms$ satisfies (CD1), because
$d>1$. Formula (CD10) with $\gamma=1$ gives exactly (CW4),
and (CD5) gives (CW5). The cited theorem proves integrability
and representative independence on the actual Borel difference
space, not just an abstract positive extension from compact
tests. Its strict-box hypotheses are precisely those retained
above. Every actual refinement remains in the original test
class used in (CW6) and in the capacity definition.
\end{proof}

\begin{corollary}[A canonical normal pair for each compact tester cluster]
\label{cor:r12-cluster-normal-pair}
For every $\nu$ as above there is a uniquely determined
variational potential $u_\nu=\Pi(\ell_\nu)\in\mathscr C$ and
normal $n_\nu=\ell_\nu-J(u_\nu)$ satisfying
\[
 \begin{gathered}
 n_\nu(w)\leq0\quad(w\in\mathscr C),\qquad n_\nu(u_\nu)=0,\\
                  e(u_\nu)\leq B_{m,M}^{d/q}.
 \end{gathered}
 \tag{CW7}
\]
The geometry-only actual probe $\psi_0$ defines a finite scalar
\[
 \Xi_\nu=\int_V(-\psi_0)\,d\widehat\mu_{u_\nu}
                       -\int_V(-\psi_0)\,d\nu\geq0,
 \tag{CW8}
\]
and, for every original compact smooth scalar $\phi$,
\[
 \left|\int_V\phi\,d(\nu-\widehat\mu_{u_\nu})\right|
                   \leq\epsilon_\phi^{-1}\Xi_\nu.
 \tag{CW9}
\]
The condition $\Xi_\nu=0$ is equivalent to existence of a
\emph{completed global zero-boundary source} for $\nu$.
It is not equivalent to compact raw-capacity no-gain.
\end{corollary}

\noindent\textbf{Status: \texttt{DIRECT\_COROLLARY}} of
(CW4), the variational normal decomposition (NI6)--(NI10),
and the original-test normal detector (NP6)--(NP10).

\begin{proof}
The preceding proposition places $\ell_\nu$ in the exact
strong-continuous dual where the variational construction is
unconditional. Since it is nonnegative on the nonpositive
potential cone, (NI9) gives
$r(u_\nu)^q=\|\ell_\nu\|_*\leq B_{m,M}$.
Raising to the power $d/q$ proves the energy bound in (CW7).
The normal relations are (NI7), and the remaining assertions
are the literal-data part of the one-probe theorem. No actual
classical representative, normalized tester bound for $u_\nu$,
or zero normal residual is inferred from variational existence.
\end{proof}

\begin{proposition}[Localization can give positive normal defect with no compact gain]
\label{prop:r12-localized-normal-defect-not-capacity-gate}
For every $q\geq1$ and every $n>q$ there are allowed original
data, one normalized actual tester $a$ and one $0\leq\chi\leq1$
as above such that the constant row $a_j=a$ has localized limit
$\nu=\chi(Da)^q\wedge S$ satisfying
\[
 \begin{gathered}
       \nu(E)\leq C_S(E;V)\quad\text{for every Borel }E\subset V,\\
       \Xi_\nu>0,
       \qquad \widehat\mu_v\ne\nu\quad\text{for every }v\in\mathscr C.
 \end{gathered}
 \tag{CW10}
\]
In particular compact no-gain holds for this row, although its
localized datum has no completed global source.
\end{proposition}

\noindent\textbf{Status: \texttt{DIRECT\_COROLLARY}}
of the proved proper-moment compact-source obstruction and
normal detector, with the constant normalized tester written
explicitly. The attempted assertion that every localized tester
cluster has $\Xi_\nu=0$ is \texttt{FALSE}.

\begin{proof}
Let $X=(\mathbf G_{m,K}^n)^{\rm an}$ over the original $K$,
$t_i=\log|T_i|$, $Q=(-1,1)^n$ and $V=t^{-1}(Q)$. Put
\[
       \beta=\sum_{i=1}^n d't_i\wedge d''t_i,
       \qquad S=\beta^{n-q}\wedge[V],
       \qquad a=\frac{1}{2n}\sum_{i=1}^n t_i^2.
 \tag{CW11}
\]
The background extends across the closed box, is original
middle-positive and separately closed, and has the required
degree. The actual smooth convex tester satisfies $0\leq a<1/2$
on the whole $V$. Its full original measure is
\[
                      (Da)^q\wedge S=n^{-q}\beta^n\wedge[V],
 \tag{CW12}
\]
which is nonzero and has positive density on the original
Gauss section by the frozen torus integration calibration.

Choose a nonzero nonnegative real smooth $\chi_0\leq1$
compactly supported in $Q$, and put $\chi=\chi_0\circ t$.
The moment $t:V\to Q$ is proper over compact subsets,
so $\chi$ is an original compact smooth scalar.
Then $\nu=\chi(Da)^q\wedge S$ is nonzero, positive,
finite and has compact tropical carrier.
The first line of (CW10) follows immediately by using the
single normalized tester $a$ in the raw capacity supremum.
No limiting or outer-regularity argument is needed.

On the other hand the proper-moment lift theorem applies
with real rank $r=n$, lift background $[V]$, and $q<n$.
It excludes every nonzero source with compact tropical
carrier from the full completed source image. Equivalently,
the smooth separator theorem supplies an original compact
smooth $b$ with
\[
       \int b\,d\nu=\nu(V)>0,\qquad
       \int b\,d\widehat\mu_v\leq0
                                      \quad(v\in\mathscr C).
 \tag{CW13}
\]
The datum is in the literal strong dual by (CW4).
The one-probe estimate applied to $b$ gives the quantitative
strictness $\Xi_\nu\geq\epsilon_b\nu(V)>0$.
This proves (CW10) on all original tests and refinements.
The special background is used solely to disprove the
universal zero-defect shortcut; the positive cluster theorem
assumes no such factorization of its arbitrary $S$.
\end{proof}

\begin{theorem}[A quantitative weakly compact family inside the strong dual]
\label{thm:r12-compact-capacity-dual-equicontinuity}
Fix a compact $H\Subset V$ and finite $m_0,M\geq0$. Let
$\mathcal G(H,m_0,M)$ be the positive Radon measures $\lambda$
with
\[
 \operatorname{supp}\lambda\subset H,\qquad
 \lambda(V)\leq m_0,\qquad
 \lambda(E)\leq M C_S^*(E;V)\quad(E\text{ Borel}).
 \tag{CW14}
\]
This is a weakly compact family, and its literal images
$\ell_\lambda$ form a weak-$*$ compact subset of
$\mathbb D_S^*$ contained entirely in
$\mathbb D_{S,\mathrm s}^*$. Here weak-$*$ means
$\sigma(\mathbb D_S^*,\mathbb D_S)$, not operator-norm
convergence.

More quantitatively, fix $R<\infty$ and let
$r(u),r(v)\leq R$, with $u,v\in\mathscr C$. Put
$\alpha=2^{-q}$ and $\varepsilon=I(u,v)^\alpha$.
There is a finite $A=A_{H,R,M,S}$, independent of $u,v$
and of the measure, such that for $0<\varepsilon\leq1$,
\[
 \sup_{\lambda\in\mathcal G(H,m_0,M)}
       \int_V|u-v|\,d\lambda
 \leq\varepsilon\left[
        m_0+\frac{2MC_0R^d}{d-1}
             +\frac{dA}{d-1}\log\frac1\varepsilon\right].
 \tag{CW15}
\]
For $I(u,v)=0$ the left side is zero. In particular this
family is uniformly continuous on every bounded-energy
strong cone, with the explicit logarithmic modulus in (CW15).

If $\lambda_j\to\lambda$ weakly in this family and
$v_j\to v$ strongly in $\mathscr C$, then
\[
 \begin{split}
   \int_V|v_j-v|\,d\lambda_j&\longrightarrow0,\\
   \int_V v_j\,d\lambda_j&\longrightarrow\int_Vv\,d\lambda.
 \end{split}
 \tag{CW16}
\]
For fixed $h\in\mathbb D_S$ one consequently has
$\ell_{\lambda_j}(h)\to\ell_\lambda(h)$.
These conclusions apply in particular to every weakly
convergent localized normalized tester row (CW2), with
$m_0$ supplied by its uniform CLN mass bound.
\end{theorem}

\noindent\textbf{Status: \texttt{NEW\_PROVED}} for the
uniform logarithmic strong modulus, the full weak-$*$
compact datum family and the jointly varying test conclusion.
Its inputs are the frozen first-energy sublevel estimate,
the all-original compact capacity estimate (EV8), completed
outer quasicontinuity and the proved datum criterion (CD).
No Banach-norm compactness or nonlinear source closure under
weak measure convergence is asserted.

\begin{proof}
The frozen compact gap estimate, passed to the strong
completion, gives a finite $C_{H,R,S}$ such that
\[
 C_S^*(H\cap\{|u-v|>t\};V)
       \leq C_{H,R,S} I(u,v)^\alpha/t\qquad(t>0).
 \tag{CW17}
\]
Here is a direct justification of this use of completion.
Choose actual defining rows for $u,v$ with uniformly
bounded energy roots, tending strongly to them. For a
fixed $t$, the limiting gap is contained in the actual
gap at $t/2$ together with the two approximation gaps
at $t/4$, outside their common null set. The latter
two capacities tend to zero. Apply (EV8) to the actual
gap and use continuity of the mixed first energies in
the strong completion. Absorb its factor two in
$C_{H,R,S}$. This uses capacity subadditivity and
capacity convergence, not continuity from above on
arbitrary compact sets or raw/outer equality.

For each $\lambda$ in (CW14), the distribution of
$|u-v|$ therefore has the simultaneous bounds
\[
 \lambda\{|u-v|>t\}
       \leq\min\left\{m_0,\frac{A\varepsilon}{t},
                            \frac{B}{t^d}\right\},
 \qquad A=MC_{H,R,S},\quad B=2MC_0R^d.
 \tag{CW18}
\]
The third bound follows from
$\{|u-v|>t\}\subset\{u<-t\}\cup\{v<-t\}$ for
nonpositive representatives, and the complete global
sublevel estimate $C_S^*\{w<-t\}\leq C_0e(w)t^{-d}$.
All pole and representative exceptional sets are
annihilated by (CW14).

Integrate (CW18) by layer cake. Use $m_0$ on
$(0,\varepsilon)$, $A\varepsilon/t$ on
$(\varepsilon,\varepsilon^{-1/(d-1)})$, and $B/t^d$
on the remaining interval. The three integrals are
respectively
\[
       m_0\varepsilon,\qquad
       \frac{dA}{d-1}\varepsilon\log(1/\varepsilon),
       \qquad \frac{B}{d-1}\varepsilon.
\]
This is (CW15). If $I=0$, the relative classes agree
outside an outer-null set, so the asserted zero follows
directly. A strongly convergent row has bounded energy
roots; hence (CW15) proves the first line of (CW16).

It remains to pass a fixed completed weight through
ordinary weak measure convergence. Put
$v^{(L)}=\max(v,-L)$. Its restriction to $H$ is bounded
and outer-quasicontinuous. For any $\delta>0$ choose
an ordinary open $G$ with $C_S^*(G;V)<\delta$ so
that this restriction is continuous on $H\setminus G$.
The ordinary continuous extension and compact cutoff
construction gives $f\in C_c(V)$ with $|f|\leq L$
and $f=v^{(L)}$ on $H\setminus G$. Thus both the
approximating measures and the limit satisfy
\[
       \int|v^{(L)}-f|\,d\lambda_j\leq2LM\delta,
       \qquad
       \int|v^{(L)}-f|\,d\lambda\leq2LM\delta.
\]
Weak convergence passes the $f$ integral; one may
equivalently use original smooth scalar density and
the common mass bound. Letting $\delta\downarrow0$
passes the fixed clipped weight. The global sublevel
bound also gives, uniformly over the entire family,
\[
 \int\bigl((-v)-L\bigr)_+\,d\lambda_j
       \leq\frac{MC_0e(v)}{d-1}L^{1-d},
\]
and the same bound for $\lambda$. Therefore
$\int v\,d\lambda_j\to\int v\,d\lambda$.
Combine this with the first line of (CW16) to obtain
its second line. Applying the fixed-weight assertion
to each of two cone members whose difference is $h$
proves convergence on all $\mathbb D_S$.

Finally, positive measures supported in the compact $H$
with mass at most $m_0$ form a weakly compact set.
The subclass (CW14) is weakly closed: on ordinary
opens use Portmanteau and
$C_S^*(G;V)=C_S(G;V)$, then take ordinary open
neighborhood infima for each Borel set, just as in
(CW6). Every member is a literal strong-dual datum
by (CD). The fixed-weight argument just proved shows
that the map $\lambda\mapsto\ell_\lambda$ is
continuous into the weak-$*$ topology. The same proof
works for nets, since its two error bounds are uniform
and its middle continuous-test limit is the definition
of weak convergence. Its image is consequently
weak-$*$ compact. This proves all the claims without
claiming that the entire strong dual is weak-$*$ closed.
\end{proof}

\paragraph{The attempted compact no-gain proof and its exact stopping point.}
The new Jensen theorem represents a nonlinear weak cluster,
before an arbitrary compact cutoff, by
$\nu=\widehat\mu_u-DV_q$ with $V_q\geq0$.
Its universal scalar loss controls every original coefficient
of $V_q$. A compact-set extremizing sequence, however, need
not minimize this Jensen loss. The finite smooth examples
in the Jensen section show that $-DV_q$ has neither sign;
convex averaging cannot be asserted to preserve, increase,
or decrease the mass on the target compact set.

Even vanishing Jensen loss yields the source of a normalized
\emph{completed} Borel potential. It does not construct a
normalized actual classical tester with at least the same
mass on that compact. Replacing a sharp compact restriction
by weak convergence is the already recorded moving-atom
failure. These are distinct obstacles, not a counterexample
to (FG1).

The variational alternative always produces the well-defined
normal pair (CW7), but (CW10) proves that its global source
defect is too strong a condition for localized compact no-gain.
Consequently the route from either scalar defect directly to
(FG1) remains a \texttt{FAILED\_ATTEMPT}. The missing
assertion is still an actual compact mass comparison for the
full tester family, not existence of a finite-energy dual
functional for each localized weak limit. No unrestricted
comparison, contact orthogonality, Radon-probe null-ideal
identification or actual-representability theorem is assumed.

\paragraph{Exact new dependencies.}
The positive dual criterion is
\path{energy/capacity_controlled_dual_data.tex}, SHA256:
\begin{quote}\small
\nolinkurl{415cab8b5265bd34d298db187317ac28cd075f4265939e0d362a0b30c42da7e0}.
\end{quote}
The variational construction is
\path{energy/nonlinear_dual_source_geometry.tex}, SHA256:
\begin{quote}\small
\nolinkurl{29a083c0f65d93f777cae9bbe0c02a36385510e7b1a6c9542552322439168743}.
\end{quote}
The normal test detector is
\path{nonlinear/single_probe_normal_residual.tex}, SHA256:
\begin{quote}\small
\nolinkurl{f09fb0f75efdce5fd637e088a70262091ed650c13cb216b5611541b6e0de28d4}.
\end{quote}
The compact-source separator is
\path{nonlinear/moment_source_separation.tex}, SHA256:
\begin{quote}\small
\nolinkurl{42c807dff24816a0e7466a089af094276fcab1a48c800821e02d28b0c098084c}.
\end{quote}
These are proofs in the present round; the frozen bounded calculus,
relative CLN and exact original torus calibration retain their
\texttt{INTERNAL\_INPUT} status.
The frozen outer gap estimate is (EV8) in
\path{R11/research_r11/comparison/first_energy_capacity_topology.tex},
SHA256:
\begin{quote}\small
\nolinkurl{5e1c0eba5cfeef5ddc15a25d4c7cb5f1c4972fc325a4568d0f01339ee203da94}.
\end{quote}
Finite local outer quasicontinuity is precisely
\texttt{prop:r11-strong-cauchy-borel-limit} in
\path{R11/research_r11/comparison/strong_first_energy_completion.tex},
SHA256:
\begin{quote}\small
\nolinkurl{11392009f63f5de3b21dd2d90d8ded48682ede8cb79665ab74e3da8ebf6534df}.
\end{quote}
